%! Author = lili
%! Date = 6/19/26

\chapter{Other possibly helpful suggestions (awaiting a better title)}\label{ch:other-possibly-helpful-suggestions-(awaiting-a-better-title)}


\section{University support}\label{sec:university-support}
May it be a laboratory space, money, or expert knowledge - it is difficult to sustain an iGEM team without
support of the institution or university the students belong to.

We are trying to stay an attractive program through press engagement (generating positive press for our university),
holding ``homecomings'' (little reception after the Jamboree) for our supporters and key people in university
policy, and through being overseen by a steering committee that mainly serves the function to help us navigate
university procedures (and thereby ensuring we do not create trouble for the university) and to manage the finances.

Though the establishment would not have been possible without our former professors Jörn Kalinowski and Alfred Pühler,
who supported and strongly advocated for iGEM\@.

Further, we established long-term advisors or instructors to minimize the loss of experience and structure through
the years and to give teams a head start organisational-wise to help them create good projects and ensure they meet
the medal criteria.
Supporters of any kind are usually more willing to help out, if the team has a track record of winning gold medals
and possibly other awards.

Of course, it is beneficial to be aligned with the teaching, scientific and technical values of your university.
Luckily, our university deeply values interdisciplinarity, which in turn leads to students from many fields being
open and interested in contributing to iGEM projects.
My favourite example for this is Michael from 2024, who studied Interdisciplinary Media Studies, and took
initiative to become part of the 2024 team.
